\section{符号说明}

本文使用的符号与单位统一列于表~\ref{tab:symbols}。全文温度用摄氏度、含水率用\textbf{干基}
（每千克干物质所含的水分质量），长度在内部计算中为米、在结果表中为厘米，时间在交付文件中为秒、在表格中为小时。
为节省篇幅，表~\ref{tab:symbols} 只列出后文反复使用的符号，一次性符号在首次出现处就地说明。

\begin{longtable}{@{}l p{7.3cm} l@{}}
  \caption{符号说明}\label{tab:symbols}\\
  \toprule
  符号 & 含义 & 单位 \\
  \midrule
  \endfirsthead
  \toprule 符号 & 含义 & 单位 \\ \midrule \endhead
  \bottomrule \endfoot
  $R_0,\ R(t),\ R_{\min}$ & 药材的初始半径、当前半径、收缩终值 & cm \\
  $L$ & 药材长度 & cm \\
  $r$ & 到药材中心的径向距离（物理坐标） & cm \\
  $\xi$ & 归一化径向坐标，等于当前物理距离与当前半径之比 & --- \\
  $t$ & 时间 & s（交付文件）/ h（表格） \\
  $T(r,t)$ & 药材内部温度 & $^\circ$C \\
  $C(r,t)$ & 药材内部干基水分浓度 & kg/kg \\
  $T_0,\ C_0$ & 温度与含水率的初值 & $^\circ$C, kg/kg \\
  $T_\infty(t),\ C_\infty(t)$ & 烘房环境温度与环境水分浓度 & $^\circ$C, kg/kg \\
  $\rho$ & 药材密度 & kg/m$^3$ \\
  $c_p$ & 药材比热容 & J/(kg$\cdot$K) \\
  $k$ & 药材导热系数 & W/(m$\cdot$K) \\
  $\alpha$ & 热扩散率，等于导热系数除以密度与比热容之积 & m$^2$/s \\
  $D$ & 水分扩散系数 & m$^2$/s \\
  $h$ & 对流传热系数 & W/(m$^2\cdot$K) \\
  $h_m$ & 对流传质系数 & m/s \\
  $\dot R$ & 药材半径的变化率 & m/s \\
  $N$ & 空间网格数（主方法取 $800$） & --- \\
  $\Delta r,\ \Delta\xi$ & 物理坐标与归一化坐标下的网格间距 & m, --- \\
  $\Delta t$ & 时间步长 & s \\
  $w_j$ & 第 $j$ 个对偶控制体的归一化度量，等于该控制体上径向坐标的积分 & --- \\
  $k_c$ & 表观对流系数，等于半径变化率与半径之比 & s$^{-1}$ \\
  $W$ & 单位长度药材内的总水量 & kg/m \\
  $C_{\text{target}}$ & 烘干判据，即全域最大干基含水率的上限 & kg/kg \\
  $t_{\text{dry}}$ & 烘干时长 & h \\
  $\text{Bi},\ \text{Bi}_m,\ \text{Fo}$ & 传热、传质 Biot 数与 Fourier 数 & --- \\
\end{longtable}
